% 拆分 -> 连接 -> 协同 -> 统一
% 4.1 拆分的艺术 拆成多图 业务和数据库拆开 RPC  2.1里面要加入业务与数据库之间的访问方式
% 业务再进行拆分
% 4.1.1 垂直拆分 按照功能切 大厅 + 战斗 大厅负责聊天 战斗则是物理计算
% 4.1.2 水平拆分 策略A 分服 这是最简单的切分 物理隔离 玩家之间不能互动 简单粗暴 不灵活 是旧时代网游 
% 策略B 分图 无缝大地图 不同服务器负责不同的坐标区域 高级扩容 技术难 体验佳 适合超级大世界

% 我们的决策 因为战斗最耗CPU 所以部署10台战斗服武器 动态副本 结合A和B 非战斗状态下 大家在一个世界 战斗时动态分配一个独立的房间 比如新开一个房间101 对应在战斗服 3的内存里new了一个Room对象

% 4.2 神经系统：RPC与服务间通信 
% 玩家A在大厅发起我要打B 这个请求是发给大厅服务器的 但是大厅服务器自己不处理战斗 它需要“指挥”战斗服去干活 这里就要引入RPC 因为大厅和战斗是两个独立的进程 它们必须通过网络函数调用协作
% 4.2.1 为什么不能同HTTP 性能差 不支持双向通信
% 4.2.2 gRPC与Protobuf
% 4.2.3 服务发现 Battle Server扩容到了10台 Gateway怎么知道它们的IP？
% 引入一个简单的注册中心

% 4.3 流量指挥官：负载均衡和网关
% 4.3.1 Nginx 只适用于登录服务器 因为它无状态 Round Robin即可 

% 4.4 最后的防线：一致性与共识
% 4.4.1 业务层的一致性 TCC事务模型 Try-Confirm-Cancel 
% 擂台
% 假设A和B战斗前要扣除10点体力作为门票 单机的时候可以不使用TCC 因为只要扣体力前用数据库的事务功能就可以保证 就算有错误也能回滚 数据库全兜底了
% 现在到了分布式阶段 网络并不可靠 Player Server 管体力 连接Redis1
% Battle Server 管房间 连接Redis2 他们之间通过RPC连接
% RPC发送前就已经扣除了玩家A的体力 但它并不知道这个体力扣除的“是否有价值” 也就是处于了RPC的未知状态 如果Battle Server开好了还好 如果没开好 这个体力白扣了
% TCC的本质就是在网络不可靠的情况下 通过“两阶段确认”来消除RPC的未知状态
% Try 冻结 
% Try1 将A的10点体力标志为预扣除的状态 同时把玩家Lock起来（类似于王者荣耀的匹配机制 进入匹配后就无法进行其他操作）
% Try2 B是否可以打架？ 是否在线 是否Lock等
% Try3 RPC询问 Battle Server的资源现在能再开一个房间吗？预留一个RoomId给我
% 所有Try结束之后才能进入Confirm 任何一步失败 直接Cancel
% Confirm 确认 这一步是幂等的 
% 正式扣除A和B的体力值 创建战斗服的房间 把A和B从匹配中改变为战斗中 
% Cancel 取消
% Try阶段有人掉链子了 那就大家恢复原状 当作什么没有发生过 
% TCC机制能够聪明的解决分布式中的不确定性的问题 业务层的目的就是保证体力不白扣  背景则是是因为数据分散在不同的物理机器上，中间隔着不可靠的网络，导致无法使用本地数据库事务，才不得不在业务代码层面模拟出一个“分布式事务”协议。


% 4.4.2 数据层的共识： Raft  -> 第五章
% 引入的注册中心 挂掉了怎么弄? 如果部署3台注册中心 怎么保证数据是一模一样的？ -> Raft
% 4.5 总结 